% !TeX root = ../../main.tex
% chapters/5_mobiwac/08_conclusion.tex -- one section of Chapter 5 (MobiWac 2026, submitted, under review).
% Split out of chapters/5_mobiwac.tex on 2026-07-28 (author request: the paper
% chapters were single long files, while the originals under articles/[mobiwac]/src/sections/
% are one file per section). The split is PURELY MECHANICAL: content was cut at
% \section boundaries and not edited. The chapter file now \inputs these in order.
% Editing rules are unchanged and still governed by the chapter's status above:
% see NORTH_STAR section 4 before changing any sentence here.
\section{Conclusion}
\label{sec:mobiwac:conclusion}

We tested whether one model can predict both the category and region of a user's next visit. The central concern was
that sharing could help one task while hurting the other. However, the evidence shows that sharing alone does not
explain the outcome because the result also depends on what the model represents. A standard place-level representation
assigns the same vector to every visit to a place, so it cannot distinguish visits made under different contexts. By
assigning one vector to each visit, the check-in-level representation can include the visit time, nearby places, and
recent visits. The controlled comparison then shows a consistent advantage for that representation on the category
task: every fold favors it at every dataset, though the size of the advantage varies from a quarter of a point to
six points. Once this information is available, the joint model can share
semantic context between the tasks while keeping a private spatial path for region prediction.

With both changes in place, joint training costs nothing that either analysis can resolve as a
meaningful loss. For region
prediction the model outperforms the dedicated model at Texas and California, and at the other four
datasets it stays within the two-point margin registered before any result was read. For category
prediction it outperforms the dedicated model at Florida; the five remaining differences are not
resolved, and every one of them is equivalent to zero within half a point. The two datasets where
the region output improves are the two with the largest region vocabularies, which is where an
auxiliary signal has the most to add, although region count and corpus size co-vary across these
states and the ordering is not monotone within the pair. The joint model also remains above every
external baseline in Table~\ref{tab:mobiwac:results} at all six datasets, by at least $3.55$ Acc@10
points over the strongest region reference (HMT-GRN, STAN, or ReHDM; the first two on our folds,
ReHDM under its own protocol) and by at least $3.06$ macro-F1 points over POI-RGNN on category.

Together, these results do not mean that multitask learning helps automatically. They show that the initial concern
about harmful sharing was incomplete: the result also depends on the representation and on how the tasks exchange
information. When the representation preserves the context of each visit and the architecture keeps a private spatial
path where the tasks differ, one model can predict both what a user will do next and where in a single forward pass.
